%!TEX program = xelatex

\documentclass[]{vvsu}

\vvsuyear{2025}

%%%%%%%%%%%%%%%%%%%

\usepackage{graphicx} % для изображений
\usepackage{tabularray} % для таблиц
\usepackage{siunitx} % для обозначений (процент, градус)
\usepackage{listings} % для листингов кода

% Список путей, где будут искаться изображения и файлы
\graphicspath{{images/}}

% Файл со списком источников (не используется)
% \addbibresource{./references.bib}

% Автор документа
\author{С.Д. Мурадуллаева}

% Настройка стилей для листингов кода
\input{listing_styles.tex}

%%%%%%%%%%%%%%%%%%%

\begin{document}

% Шапка
\vvsuhead{\linespread{1}\selectfont{}МИНОБРНАУКИ РОССИИ\\
\vspace{10pt}Федеральное государственное бюджетное образовательное учреждение\\
высшего образования\\
\fontsize{13}{13}\selectfont{}<<ВЛАДИВОСТОКСКИЙ ГОСУДАРСТВЕННЫЙ УНИВЕРСИТЕТ>>\\
(ФГБОУ ВО <<ВВГУ>>)\\
\vspace{10pt}\fontsize{12}{12}\selectfont{}ИНСТИТУТ ИНФОРМАЦИОННЫХ ТЕХНОЛОГИЙ И АНАЛИЗА ДАННЫХ\\
КАФЕДРА ИНФОРМАЦИОННЫХ ТЕХНОЛОГИЙ И СИСТЕМ}

% Название отчета
\title{Отчет\\по лабораторной работе №4}
\subtitle{по дисциплине\\<<Информатика и основы программирования>>}

% Участники работы
\member{Студент\\ гр. БИН-25-3}{С.Д. Мурадуллаева}
\member{Ассистент\\ преподавателя}{М.В. Водяницкий}

% Вывод титульника
\maketitle

% Задание
% Задание
\begin{addition}{Задание}
  Выполнить задания на Python и оформить отчет по стандартам ВВГУ.

  \textit{\textbf{Задание 1.}}  
  Написать программу, которая определяет, как будет вести себя кондиционер. Если температура в помещении 20 градусов и выше, то кондиционер выключается, если меньше - включается. Температура должна вводится пользователем с консоли.

  Пример:
  \begin{verbatim}
  Введите температуру: 18
  Кондиционер включен\end{verbatim}

  \textit{\textbf{Задание 2.}}  
  Год делится на четыре сезона: зима, весна, лето и осень. Написать программу, которая запрашивает у пользователя номер месяца и выводит к какому сезону этот месяц относится.

  Пример:
  \begin{verbatim}
  Введите номер месяца: 4
  Это весна\end{verbatim}
  
  \textit{\textbf{Задание 3.}}  
  Считается, что один год, прожитый собакой, эквивалентен семи человеческим годам. При этом зачастую не учитывается, что собаки становятся абсолютно взрослыми уже к двум годам. Таким образом, многие предпочитают каждый из первых двух лет жизни собаки приравнивать к 10.5 годам человеческой жизни, а все последующие к 4.
  Написать программу, которая будет переводить собачий возраст в человеческий. Программа должна корректно обрабатывать входные данные и выводить соответствующие сообщения об ошибках:
  \begin{vvsu_list}
    \item Если вводится не число
    \item Если вводится число меньше 1
    \item Если вводится число большее 22
  \end{vvsu_list}
  Пример:
  \begin{verbatim}
  Введите возраст собаки (в годах): 5
  Возраст собаки в человеческих годах: 33.0\end{verbatim}
  Пример:
  \begin{verbatim}
  Введите возраст собаки (в годах): 0  
  Ошибка: возраст должен быть не меньше 1\end{verbatim}

  \textit{\textbf{Задание 4.}}  
  Число делиться на 6 только в случае соблюдения двух условий:
  \begin{vvsu_list}
    \item Последняя цифра четная
    \item Сумма всех цифр делиться на 3
  \end{vvsu_list}
  Написать программу, которая выведет делиться ли введенное число на 6 или нет.

  \textit{\textbf{Задание 5.}}  
  Написать программу, которая будет проверять пароль на надежность. Пароль считается надежным, если его длина не менее 8 символов и если он содержит:
  \begin{vvsu_list}
    \item Заглавные буквы латиницы
    \item Строчные буквы латиницы
    \item Числа
    \item Специальные знаки
  \end{vvsu_list}
  В случае, если пароль не проходит по одному из условий, необходимо сообщить пользователю каким именно условиям он не удовлетворяет.

  Пример:
  \begin{verbatim}
  Введите пароль: qwerty  
  Пароль ненадежный: отсутствуют заглавные буквы, числа и специальные символы\end{verbatim}
  
  \textit{\textbf{Задание 6.}}  
  Написать программу, которая определяет, является ли введенный пользователем год високосным. Год считается високосным, если он делится на 4, но не делится на 100, либо если он делится на 400.

  Пример:
  \begin{verbatim}
  Введите год: 2024  
  2024 - високосный год\end{verbatim}

  \textit{\textbf{Задание 7.}}  
  Написать программу, которая запрашивает у пользователя три числа и выводит на экран наименьшее из них. При решении нельзя использовать встроенные функции min() и max().

  Пример:
  \begin{verbatim}
  Введите три числа: 8 3 5  
  Наименьшее число: 3\end{verbatim}

  \textit{\textbf{Задание 8.}}  
  В магазине проводится акция. Акция работает по следующим правилам:

  \begin{vvsu_table}{Сумма покупки и скидка}
    \label{table:sum_sale}
    \begin{tblr}{
      width=1\linewidth,
      colspec={lX},
      hlines, vlines
    }
    Сумма покупки   & Скидка \\
    до 1000    & 0\% \\
    1000–5000 & 5\% \\
    5000–10000      & 10\% \\
    более 10000      & 15\% \\
    \end{tblr}
  \end{vvsu_table}

  Напишите программу, которая запрашивает сумму покупки и выводит размер скидки и итоговую сумму к оплате.

  Пример:
  \begin{verbatim}
  Введите сумму покупки: 7500  
  Ваша скидка: 10%  
  К оплате: 6750.0\end{verbatim}

  \textit{\textbf{Задание 9.}}  
  Написать программу, которая определяет время суток по введенному часу (целое число от 0 до 23).

  \begin{vvsu_table}{Время и период}
    \label{table:time_period}
    \begin{tblr}{
      width=1\linewidth,
      colspec={lX},
      hlines, vlines
    }
    Время   & Период \\
    0–5    & Ночь \\
    6–11 & Утро \\
    12–17      & День \\
    18–23      & Вечер \\
    \end{tblr}
  \end{vvsu_table}

  Пример:
  \begin{verbatim}
  Введите час (0–23): 20  
  Сейчас вечер\end{verbatim}

   Напишите программу, которая запрашивает сумму покупки и выводит размер скидки и итоговую сумму к оплате.

  Пример:
  \begin{verbatim}
  Введите сумму покупки: 7500  
  Ваша скидка: 10%  
  К оплате: 6750.0\end{verbatim}

  \textit{\textbf{Задание 9.}}  
  Написать программу, которая определяет время суток по введенному часу (целое число от 0 до 23).

  \begin{vvsu_table}{Время и период}
    \label{table:time_period}
    \begin{tblr}{
      width=1\linewidth,
      colspec={lX},
      hlines, vlines
    }
    Время   & Период \\
    0–5    & Ночь \\
    6–11 & Утро \\
    12–17      & День \\
    18–23      & Вечер \\
    \end{tblr}
  \end{vvsu_table}

  Пример:
  \begin{verbatim}
  Введите час (0–23): 20  
  Сейчас вечер\end{verbatim}

 % Подглава - Задание 1
\subsection{Задание 1} 
В данном задании необходимо написать программу, которая определяет включен или выключен кондиционер.
Пользователь вводит температуру. Далее в зависимости от введенной пользователем температуры определяется выключен кондиционер или включен.
На рисунке \ref{fig:code_task_1} представлен код полученной программы.

\begin{vvsu_figure}{Листинг программы для задания 1}{fig:code_task_1}
  \begin{minipage}{.75\textwidth}
    \lstinputlisting[language=Python,basicstyle=\fontsize{10}{10}\linespread{1}\selectfont\ttfamily]{code/task1.py}
  \end{minipage}
\end{vvsu_figure}

Пояснение работы программы:
\begin{vvsu_list}
  \item просит ввести температуру и превращает текст в число
  \item проверяет, если температура 20 или больше
  \item если условия верно, выключает кондиционер
  \item если условие не верно (температура меньше 20)
  \item включает кондиционер
\end{vvsu_list}

% Подглава - Задание 2
\subsection{Задание 2}

В данном задании необходимо определить время года по номеру месяца.
Пользователь вводит номер месяца. Далее в зависимости от введенного пользователем номера месяца определяется время года.
На рисунке унке \ref{fig:code_task_2} представлен код программы.

\begin{vvsu_figure}{Листинг программы для задания 2}{fig:code_task_2}
  \begin{minipage}{.75\textwidth}
    \lstinputlisting[language=Python,basicstyle=\fontsize{10}{10}\linespread{1}\selectfont\ttfamily]{code/task2.py}
  \end{minipage}
\end{vvsu_figure}

Пояснение работы программы:
\begin{vvsu_list}
  \item надо ввести номер месяца 
  \item если мы введем цифры 12 1 2
  \item мы получим сообщение 'это зима'
  \item если мы введем цифры 3 4 5
  \item мы получим сообщение 'это весна'
  \item если мы введем цифры 6 7 8
  \item мы получим сообщение 'это лето'
  \item если мы введем цифры 9 10 11
  \item мы получим сообщение 'это осень'
  \item else: срабатывает если все условия выше не подошли (месяц не от 1 до 12)
\end{vvsu_list}

% Подглава - Задание 3
\subsection{Задание 3}

В данном задании необходимо написать программу, которая переводит возраст собаки в человеческие годы.
Пользователь вводит возраст собаки. Проверяется соответвует ли введенное число условию, и только потом проиводится перевод возраста собаки в человеческие года. Полученный результат выводится в консоль. 
На рисунке \ref{fig:code_task_3} представлен код программы.

\begin{vvsu_figure}{Листинг программы для задания 3}{fig:code_task_3}
  \begin{minipage}{.75\textwidth}
    \lstinputlisting[language=Python,basicstyle=\fontsize{10}{10}\linespread{1}\selectfont\ttfamily]{code/task3.py}
  \end{minipage}
\end{vvsu_figure}

Пояснение работы программы:
\begin{vvsu_list}
  \item С помощью функции input() пользователем присваивается значение переменной возраста собаки.
  \item Создается функция, в которой через try except проверятся является ли введенная переменная числом.
  \item Если данная функция является числом, то возвращается True, а если нет, то возвращается False.
  \item С помощью условного оператора if проверятся, что возвращается.
  \item Если возвращается False, то в консоль выводится сообщение о соответвующей ошибке.
  \item Если возвращается True, то переменная возраста собаки переводится из str в int, с помощью функции int().
  \item Через условный оператор if проверятеся меньше ли переменная возраста собаки, чем 1.
  \item Если переменная возраста собаки меньше, чем 1, то в консоль выводится соответствующее сообщение об ошибке.
  \item Через условный оператор if проверятеся больше ли переменная возраста собаки, чем 20.
  \item Если переменная возраста собаки больше, чем 20, то в консоль выводится соответствующее сообщение об ошибке.
  \item Через условный оператор if проверятеся равна ли переменная возраста собаки 1.
  \item Если переменная возраста собаки равна 1, то выводится сообщение о том, что возраст собаки в человеческих годах равен 10.5.
  \item Через условный оператор if проверятеся равна ли переменная возраста собаки 2.
  \item Если переменная возраста собаки равна 2, то выводится сообщение о том, что возраст собаки в человеческих годах равен 21.
  \item Через условный оператор if проверяется, больше ли переменная возраста собаки, чем 2.
  \item Если переменная возраста собаки больше 2, то с помощью цикла for высчитывается возраст собаки в человеческих годах и соответствующий результат выводится в консоль.
\end{vvsu_list}

% Подглава - Задание 4
\subsection{Задание 4}

Данная программа определяет делится ли введенное пользователем с консоли число на 6.
На рисунке \ref{fig:code_task_4} представлен код решения.

\begin{vvsu_figure}{Листинг программы для задания 4}{fig:code_task_4}
  \begin{minipage}{.75\textwidth}
    \lstinputlisting[language=Python,basicstyle=\fontsize{10}{10}\linespread{1}\selectfont\ttfamily]{code/task4.py}
  \end{minipage}
\end{vvsu_figure}

Пояснение работы программы:
\begin{vvsu_list}
  \item Создается переменная число, значение которой вводится пользователем, с помощью функции input().
  \item Создается массив, который состоит из цифр заданного числа.
  \item С помощью условного оператора if проверяются два условия: четная ли последняя цифра и делится ли сумма всех цифр на 3.
  \item Если условия выполняются, в консоль выводится сообщение о том, что число делится на 6.
  \item Если условия не выполняются, в консоль выводится сообщение о том, что число не делится на 6.
\end{vvsu_list}

После выполнения программы в консоли отображается вычисленное\linebreak значение выражения.

% Подглава - Задание 5
\subsection{Задание 5}

В данном задании требуется написать программу, которая будет проверять пароль на надежность по заданным условиям.
Пользователь вводит пароль. Введенный пароль проверяется на надежность и выводится соответствующее сообщение.
На рисунке \ref{fig:code_task_5} представлен код программы.

\begin{vvsu_figure}{Листинг программы для задания 5}{fig:code_task_5}
  \begin{minipage}{.75\textwidth}
    \lstinputlisting[language=Python,basicstyle=\fontsize{10}{10}\linespread{1}\selectfont\ttfamily]{code/task5.py}
  \end{minipage}
\end{vvsu_figure}

Пояснение работы программы:
\begin{vvsu_list}
  \item import string- импортирует модуль для работы со строками
\end{vvsu_list}

% Подглава - Задание 6
\subsection{Задание 6}
 
В данном задании требуется написать программу, которая определяет, является ли год високосным.
Пользователем вводится год. Заданный год проверяется по условиям и в консоль выводится соответствующее сообщение.
На рисунке \ref{fig:code_task_6} представлен код программы.

\begin{vvsu_figure}{Листинг программы для задания 6}{fig:code_task_6}
  \begin{minipage}{.75\textwidth}
    \lstinputlisting[language=Python,basicstyle=\fontsize{10}{10}\linespread{1}\selectfont\ttfamily]{code/task6.py}
  \end{minipage}
\end{vvsu_figure}

Пояснение работы программы:
\begin{vvsu_list}
  \item Создается переменная года, значение которой задается пользователем, с помощью функции input().
  \item Тип данных заданной переменной переводится из str в year, с помощью функции year().
  \item С помощью условного оператора if переменная года проверяется по двум условиям: делится ли значение заданной переменной на 4 и при этом не делится на 100, и делится ли значение заданной переменной на 400.
  \item Если заданная переменная соответвует данным условиям, то в консоль выводится сообщение от том, что заданный год является високосным.
  \item Если заданная переменная не соответвует данным условиям, то в консоль выводится сообщение от том, что заданный год не является високосным.
\end{vvsu_list}

% Подглава - Задание 7
\subsection{Задание 7}

В данном задании необходимо написать программу, которая запрашивает у пользователя три числа и выводим на экран наименьшее из них. Использовать встроенные функции min() и max() при решении данного задания запрещается.
Создается переменная, в которую пользователем вводятся три числа через пробел. Создается массив, содержащий заданные три числа и с помощью сортировки выводится наименьшее число.  
На рисунке \ref{fig:code_task_7} представлен код программы.

\begin{vvsu_figure}{Листинг программы для задания 7}{fig:code_task_7}
  \begin{minipage}{.75\textwidth}
    \lstinputlisting[language=Python,basicstyle=\fontsize{10}{10}\linespread{1}\selectfont\ttfamily]{code/task7.py}
  \end{minipage}
\end{vvsu_figure}

Пояснение работы программы:
\begin{vvsu_list}
  \item Создается переменная, куда пользователь вводит три числа через пробел, с помощью функции input().
  \item Создается массив из чисел заданной переменной, с помощью метода .split().
  \item Массив сортируется от меньшего к большему с помощью функции sorted().
  \item Первая переменная, содержащаяся в массиве, которая и является наименьшим числом, выводится в консоль.
\end{vvsu_list}

После выполнения программы в консоли выводятся результаты всех арифметических и логических операций между переменными.

% Подглава - Задание 8
\subsection{Задание 8}

В данном задании нужно написать программу, которая запрашивает сумму покупки и выводит размер скидки и итоговую сумму к оплате.  
Пользователь вводит сумму покупок. В зависимости из размеры суммы покупок вычисляется скидка. Размер скидки и итоговая сумма к оплате выводится в консоль.
На рисунке \ref{fig:code_task_8} представлен код программы.

\begin{vvsu_figure}{Листинг программы для задания 8}{fig:code_task_8}
  \begin{minipage}{.75\textwidth}
    \lstinputlisting[language=Python,basicstyle=\fontsize{10}{10}\linespread{1}\selectfont\ttfamily]{code/task8.py}
  \end{minipage}
\end{vvsu_figure}

Пояснение работы программы:
\begin{vvsu_list}
  \item Создается переменная суммы покупок, значение которой присваивается пользователем с помощью функции input().
  \item Тип данных заданной переменной переводится из str в float.
  \item В строчках 2-17, с помощью условного оператора if, проверяются условия.
  \item В зависимости от выполненного условия вычисляются размер скидки и итоговая сумма к оплате.
  \item Размер скидки и итоговая сумма к оплате выводятся в консоль.
\end{vvsu_list}

% Подглава - Задание 9
\subsection{Задание 9}

В данном задании необходимо написать программу, которая определяет время суток по введенному часу.  
Пользователем вводится час. В зависимости от введенного часа, в консоль выводится время суток. 
На рисунке \ref{fig:code_task_9} представлен код программы.

\begin{vvsu_figure}{Листинг программы для задания 9}{fig:code_task_9}
  \begin{minipage}{.75\textwidth}
    \lstinputlisting[language=Python,basicstyle=\fontsize{10}{10}\linespread{1}\selectfont\ttfamily]{code/task9.py}
  \end{minipage}
\end{vvsu_figure}

Пояснение работы программы:
\begin{vvsu_list}
  \item Создается переменная часа, значение которой присваивается пользователем, с помощью функции input().
  \item Тип данных заданной переменной переводится из str в int.
  \item С помощью условного оператора if переменная часа проверяется по заданным условиям.
  \item Если заданная переменная часа меньше 6, то в консоль выводится сообщение о том, что сейчас ночь.
  \item Если заданная переменная часа меньше 11, но больше или равна 6, то в консоль выводится сообщение о том, что сейчас утро.
  \item Если заданная переменная часа меньше 17, но больше или равна 12, то в консоль выводится сообщение о том, что сейчас день.
  \item Если заданная переменная часа меньше или равна 23, но больше или равна 18, то в консоль выводится сообщение о том, что сейчас вечер.
\end{vvsu_list}

Результатом выполнения программы является восстановленное предложение:
<<Съешь ещё этих мягких французских булок, да выпей чаю>>.

% Подглава - Задание 10
\subsection{Задание 10}

В данном задании необходимо написать программу, которая определяет, является ли число простым. Программа должна корректно обрабатывать некорректный ввод и выводить соответствующее сообщения об ошибках.  
Пользователем вводится что-то. Если было введено не число, то выводится соответствующее сообщение об ошибке и программа завершается, если было введено число, то - продолжается. Далее проверяется, простое число или нет, и в консоль выводится соответствующее сообщение.
На рисунке \ref{fig:code_task_10} представлен код программы.

\begin{vvsu_figure}{Листинг программы для задания 10}{fig:code_task_10}
  \begin{minipage}{.75\textwidth}
    \lstinputlisting[language=Python,basicstyle=\fontsize{10}{10}\linespread{1}\selectfont\ttfamily]{code/task10.py}
  \end{minipage}
\end{vvsu_figure}

Пояснение работы программы:
\begin{vvsu_list}
  \item Создается переменная, значение которой присваивается пользователем.
  \item Создается переменная check, в которой проверяется является заданная переменная числом или нет, с помощью метода .isdigit(). Переменной присваивается значение True или False соответственно.
  \item Если переменная check содержит значение False, то в консоль выводится соответствующее сообщение об ошибке и программа завершается.
  \item Создается переменная isprime, которая содержит значение True.
  \item Если переменная check содержит значение True, то тип заданной переменной переводится из str в int, с помощью функции int().
  \item Заданная переменная проверяется по условию и если она по нему подходит, то значение переменной isprime меняется с True на False.
  \item С помощью условного оператора if проверяются два условия: больше ли заданная переменная, чем 1 и равна ли переменная isprime True.
  \item Если условия выполняются, то выводится сообщение о том, что заданное число - простое.
  \item С помощью условного оператора if проверяются два условия: больше ли заданная переменная, чем 1 и равна ли переменная isprime False.
  \item Если условия выполняются, то выводится сообщение о том, что заданное число - составное.
\end{vvsu_list}

% \begin{application}{Приложение А}{Приложение А\\ \vspace{1em} Описание процесса работы с драйверами}
%   Самое обычное тестовое приложение
% \end{application}

\end{document}